\section{条件信息、核与贝叶斯公式}\label{sec:05-02}

\subsection{条件期望：Radon--Nikodym 的信息版本}\label{subsec:5-2-1}
\index{条件期望}\index{版本}\index{Doob--Dynkin 因子化}

设 $(\Omega,\mathcal F,\Prob)$ 是概率空间。给定 $X\in L^1(\Prob)$，若观察者能看见
$\mathcal F$ 中的全部事件，那么 $X$ 本身就是可用资料；若观察者只能分辨一个子
$\sigma$-代数 $\mathcal G\subset\mathcal F$ 中的事件，我们就必须把 $X$ 压缩成一个
$\mathcal G$-可测变量。压缩不能随意进行：在每个可观察事件上，它应当保留 $X$ 的总平均。

先看一个没有零概率事件搅局的模型。

\begin{example}[只报奇偶时如何预测赔付]\label{ex:5-2-1-prediction-pressure}
公平骰子的点数记为 $D$，赔付为 $X=D^2$。理赔员只收到“奇数”或“偶数”这一位资料，故可用的
预测必为
\[
 Y=a\1_{\{D\text{ 为奇数}\}}+b\1_{\{D\text{ 为偶数}\}}.
\]
其均方误差为
\[
 \frac16\sum_{d\in\{1,3,5\}}(d^2-a)^2
 +\frac16\sum_{d\in\{2,4,6\}}(d^2-b)^2.
\]
两组平方和可逐项展开为
\begin{align*}
 \sum_{d\in\{1,3,5\}}(d^2-a)^2
 &=3a^2-70a+707
 =3\left(a-\frac{35}{3}\right)^2+\frac{896}{3},\\
 \sum_{d\in\{2,4,6\}}(d^2-b)^2
 &=3b^2-112b+1568
 =3\left(b-\frac{56}{3}\right)^2+\frac{1568}{3}.
\end{align*}
两个平方项的系数都严格为正，因此唯一极小点为
\[
 a=\frac{1+9+25}{3}=\frac{35}{3},
 \qquad
 b=\frac{4+16+36}{3}=\frac{56}{3}.
\]
这个答案还有一个不涉及微分的刻画。若 $A_o=\{D\text{ 为奇数}\}$、
$A_e=\{D\text{ 为偶数}\}$，则
\[
 \E[Y\1_{A_o}]=\frac{35}{6}=\E[X\1_{A_o}],
 \qquad
 \E[Y\1_{A_e}]=\frac{28}{3}=\E[X\1_{A_e}].
\]
具体地，在奇数块上
\[
 \E[Y\1_{A_o}]=\frac12\cdot\frac{35}{3}=\frac{35}{6},
 \qquad
 \E[X\1_{A_o}]=\frac{1^2+3^2+5^2}{6}=\frac{35}{6},
\]
偶数块上两边同为
$(4+16+36)/6=28/3$。也就是说，预测在每个已知信息块上保留了原变量的积分。后一种刻画不需要把信息分割列成有限张
清单，因而可以推广到连续观测。
\end{example}

初等概率中常见的比值
$\E[X\mid Z=y]$ 在 $\Prob(Z=y)>0$ 时可以计算，但连续分布通常满足
$\Prob(Z=y)=0$。在零分母前继续写比值不会产生定义。正确的次序是先给定可观察的
$\sigma$-代数并保留其上全部积分，再在状态空间足够良好时把所得对象组织成关于 $y$ 的核。

\begin{definition}[条件期望]\label{def:5-2-1-1}
设 $X\in L^1(\Prob)$ 且 $\mathcal G\subset\mathcal F$ 是子 $\sigma$-代数。若随机变量
$Y$ 满足
\begin{enumerate}
  \item $Y$ 是 $\mathcal G$-可测且 $Y\in L^1(\Prob)$；
  \item 对每个 $G\in\mathcal G$，
  \[
    \int_GY\,\dd\Prob=\int_GX\,\dd\Prob,
  \]
\end{enumerate}
则称 $Y$ 是 $X$ 关于 $\mathcal G$ 的一个条件期望，记作
$Y=\E[X\mid\mathcal G]$。
\end{definition}

\begin{definition}[关于随机变量的条件期望]\label{def:5-2-1-2}
若 $Z$ 是随机元，则
\[
 \E[X\mid Z]:=\E[X\mid\sigma(Z)].
\]
这是关于 $\sigma(Z)$ 的等价类记号，而不是先对每个 $z$ 定义一个比值。若 $Z$ 取值于标准
Borel 空间，\cref{proposition:5-2-1-doob-dynkin} 将给出一个 Borel 函数 $g$，使可选择
$\E[X\mid Z]=g(Z)$。
\end{definition}

\begin{definition}[版本]\label{def:5-2-1-3}
两个几乎处处相等的 $\mathcal G$-可测函数代表同一个条件期望等价类；其中任一具体代表称为一个
\emph{版本}。版本只在 $\Prob$-a.s. 意义下唯一。若等式带有不可数多个时间或参数，逐参数成立的
零测例外集未必能合并成一个零测集，因而还需另行选择相容版本。
\end{definition}

第~3~章已经从 Radon--Nikodym 定理得到过相同的积分表示。现在把它正式识别为条件期望。

\begin{theorem}[条件期望的存在唯一性]\label{theorem:5-2-1-1}
对每个 $X\in L^1(\Prob)$ 与每个子 $\sigma$-代数
$\mathcal G\subset\mathcal F$，条件期望 $\E[X\mid\mathcal G]$ 存在，并在
$\Prob$-a.s. 意义下唯一。
\end{theorem}

\begin{proof}
先设 $X\geq0$。在 $(\Omega,\mathcal G)$ 上定义
\[
 \nu(G)=\int_GX\,\dd\Prob.
\]
因为 $X\in L^1$，$\nu$ 是有限正测度；若 $\Prob(G)=0$，则 $\nu(G)=0$，故
$\nu\ll\Prob|_{\mathcal G}$。由 Radon--Nikodym 定理
\ref{theorem:3-3-1-1}，存在 $\mathcal G$-可测 $Y\geq0$ 使
\[
 \nu(G)=\int_GY\,\dd\Prob\qquad(G\in\mathcal G).
\]
取 $G=\Omega$ 得 $\int Y=\int X<\infty$，所以 $Y\in L^1$。

若 $X$ 为实值，分别对 $X^+$ 与 $X^-$ 作上述构造，记所得密度为 $Y^+,Y^-$。二者均可积，
故 $Y=Y^+-Y^-$ 可积，并在每个 $G\in\mathcal G$ 上满足所需积分等式。若 $X$ 为复值，分别
处理 $\Re X$ 与 $\Im X$，再组成 $Y=Y_1+iY_2$。

最后证明唯一性。设实值 $Y_1,Y_2$ 都满足定义，并令
$A=\{Y_1>Y_2\}\in\mathcal G$。两条积分保持式给出
\[
 \int_A(Y_1-Y_2)\,\dd\Prob=0.
\]
被积函数在 $A$ 上严格为正；非负函数积分为零当且仅当它几乎处处为零，故
$\Prob(A)=0$。交换 $Y_1,Y_2$ 得 $\Prob(Y_2>Y_1)=0$，所以二者 a.s. 相等。复值情形对
实部和虚部分别使用这个结论。
\end{proof}

“关于 $Z$ 可测”还没有自动变成“是 $Z$ 的函数”。下面的证明把这一点落实到一个可数水平集族，
从而避免为每个水平集分别选像后直接跳到因子化。

\begin{proposition}[Doob--Dynkin 因子化]\label{proposition:5-2-1-doob-dynkin}
设 $Z:\Omega\to T$ 取值于标准 Borel 空间，实值随机变量 $Y$ 对 $\sigma(Z)$ 可测。则存在
Borel 函数 $g:T\to\R$ 使 $Y=g(Z)$。特别地，对 $X\in L^1$ 可选择一个 Borel 函数 $g$，
使 $\E[X\mid Z]=g(Z)$ a.s.
\end{proposition}

\begin{proof}
对每个 $q\in\Q$，集合 $\{Y<q\}$ 属于 $\sigma(Z)$，故可取 Borel 集
$A_q\subset T$ 使
\[
 Z^{-1}(A_q)=\{Y<q\}.
\]
这些 $A_q$ 未必彼此相容。令
\[
 B_q=\bigcup_{r<q,\ r\in\Q}A_r.
\]
于是 $B_q$ 随 $q$ 非降，并且
\[
 B_q=\bigcup_{s<q,\ s\in\Q}B_s,
 \qquad
 Z^{-1}(B_q)=\bigcup_{r<q}\{Y<r\}=\{Y<q\}.
\]
两个集合
\[
 N_+=T\setminus\bigcup_{q\in\Q}B_q,
 \qquad
 N_- =\bigcap_{q\in\Q}B_q
\]
都是 Borel 集，而且都不与 $Z(\Omega)$ 相交：给定 $\omega$，总能取有理数
$q>Y(\omega)$ 与 $r<Y(\omega)$。先定义扩展实值函数
\[
 h(t)=\inf\{q\in\Q:t\in B_q\},
\]
其中空集的下确界取 $+\infty$。对每个 $a\in\R$，
\[
 \{h<a\}=\bigcup_{q<a,\ q\in\Q}B_q,
\]
故 $h$ 可测。令 $g=h$ 于 $T\setminus(N_+\cup N_-)$，并在
$N_+\cup N_-$ 上令 $g=0$，则 $g$ 是处处实值的 Borel 函数。对每个 $\omega$，
\[
 \{q\in\Q:Z(\omega)\in B_q\}=\{q\in\Q:Y(\omega)<q\},
\]
这个集合的下确界等于 $Y(\omega)$，所以 $g(Z)=Y$。

将第一部分用于条件期望的任一实值版本即可。复值条件期望对实部、虚部分别因子化。
\end{proof}

条件期望还需要容纳非负但不可积的变量；稍后核表示中取任意非负测试函数时会用到这一扩展。

\begin{proposition}[非负变量的扩展条件期望]\label{proposition:5-2-1-nonnegative}
设 $X:\Omega\to[0,\infty]$ 可测，不要求 $X$ 可积。可定义一个
$[0,\infty]$ 值 $\mathcal G$-可测变量 $\E[X\mid\mathcal G]$，使
\begin{equation}\label{eq:5-2-extended-ce}
 \int_G\E[X\mid\mathcal G]\,\dd\Prob
 =\int_GX\,\dd\Prob\qquad(G\in\mathcal G),
\end{equation}
两边允许为 $+\infty$。它在 a.s. 意义下唯一。若 $0\leq X_n\uparrow X$，则可相容地选择版本使
\begin{equation}\label{eq:5-2-conditional-mct}
 \E[X_n\mid\mathcal G]\uparrow\E[X\mid\mathcal G]
 \quad\text{a.s.}
\end{equation}
\end{proposition}

\begin{proof}
先记录可积变量的单调性。若 $U,V\in L^1$ 且 $U\leq V$ a.s.，取条件期望版本
$U_0,V_0$，并令 $A=\{U_0>V_0\}\in\mathcal G$。则
\[
 \int_A(U_0-V_0)\,\dd\Prob
 =\int_A(U-V)\,\dd\Prob\leq0.
\]
左端是非负函数的积分，故 $\Prob(A)=0$。所以 $U_0\leq V_0$ a.s.

令 $X_n=X\wedge n$。逐个选取 $Y_n=\E[X_n\mid\mathcal G]$ 的非负版本。由刚证的单调性，
$Y_n\leq Y_{n+1}$ a.s.；删除这些例外集的可数并，并在该并集上把所有 $Y_n$ 改为零，便可假设
$Y_n(\omega)$ 对每个 $\omega$ 都非降。定义 $Y=\lim_nY_n$。它非负且 $\mathcal G$-可测。
对每个 $G\in\mathcal G$，两次使用单调收敛定理得到
\[
 \int_GY\,\dd\Prob
 =\lim_n\int_GY_n\,\dd\Prob
 =\lim_n\int_G(X\wedge n)\,\dd\Prob
 =\int_GX\,\dd\Prob.
\]
这证明 \eqref{eq:5-2-extended-ce}。

若 $Y,Z$ 都满足 \eqref{eq:5-2-extended-ce}，对 $m,n\geq1$ 令
\[
 A_{m,n}=\{Y\geq Z+1/m\}\cap\{Z\leq n\}\in\mathcal G.
\]
在 $A_{m,n}$ 上 $Z$ 的积分至多为 $n$，故两条保持式说明 $Y,Z$ 在该集合上的积分都是同一个
有限数。于是
\[
 0=\int_{A_{m,n}}(Y-Z)\,\dd\Prob
 \geq \frac1m\Prob(A_{m,n}),
\]
所以 $A_{m,n}$ 为零测集。对 $m,n$ 取可数并得 $\Prob(Y>Z)=0$；交换二者便得唯一性。

最后，若 $0\leq X_n\uparrow X$，则对每个 $n,m$ 都有
$X_n\wedge m\leq X_{n+1}\wedge m$。对这些可积截断使用已证的单调性，再合并关于
$(n,m)$ 的可数个例外集并令 $m\to\infty$，得到
\[
 \E[X_n\mid\mathcal G]\leq\E[X_{n+1}\mid\mathcal G]\quad\text{a.s.}
\]
在一个共同零集上修改版本后，令 $W=\lim_n\E[X_n\mid\mathcal G]$。对每个
$G\in\mathcal G$，单调收敛给
$\int_GW=\lim_n\int_GX_n=\int_GX$。扩展条件期望的唯一性说明
$W=\E[X\mid\mathcal G]$ a.s.，即 \eqref{eq:5-2-conditional-mct}。
\end{proof}

\begin{example}[不可积变量的条件平均可以为无穷]\label{ex:5-2-1-nonintegrable}
在 $((0,1),\Borel((0,1)),m)$ 上令 $X(\omega)=\omega^{-1}$，并令
$\mathcal G=\{\varnothing,(0,1)\}$。每个 $\mathcal G$-可测非负函数都是常数，且
\[
 \E[X\wedge n]
 =\int_0^{1/n}n\,\dd\omega+\int_{1/n}^1\frac{\dd\omega}{\omega}
 =1+\log n\longrightarrow+\infty.
\]
所以 $\E[X\mid\mathcal G]$ 是常数 $+\infty$。扩展条件期望并不承诺有限值；它承诺的是
非负积分恒等式。
\end{example}

下面几条运算都从积分刻画推出，而不是从点态条件概率推出。

\begin{proposition}[基本运算与 $L^1$ 收缩]\label{proposition:5-2-1-2}
对 $X,X_1,X_2\in L^1$、标量 $a,b$，条件期望满足
\[
 \E[aX_1+bX_2\mid\mathcal G]
 =a\E[X_1\mid\mathcal G]+b\E[X_2\mid\mathcal G].
\]
若 $X\geq0$，则 $\E[X\mid\mathcal G]\geq0$ a.s.；并且对实值或复值 $X$，
\begin{equation}\label{eq:5-2-l1-contraction}
 \bigl|\E[X\mid\mathcal G]\bigr|
 \leq\E[|X|\mid\mathcal G]\quad\text{a.s.},
 \qquad
 \|\E[X\mid\mathcal G]\|_1\leq\|X\|_1.
\end{equation}
\end{proposition}

\begin{proof}
线性组合是 $\mathcal G$-可测可积函数，并在每个 $G\in\mathcal G$ 上具有正确积分；唯一性给出
线性。正性已包含在存在性证明对非负 Radon--Nikodym 密度的选择中，也可由上一证明的单调性取
$U=0,V=X$ 得到。

若 $X$ 为实值，则 $-|X|\leq X\leq|X|$。单调性和线性给
\[
 -\E[|X|\mid\mathcal G]
 \leq\E[X\mid\mathcal G]
 \leq\E[|X|\mid\mathcal G],
\]
即第一条模长不等式。

若 $X$ 为复值，记 $Y=\E[X\mid\mathcal G]$、
$Z=\E[|X|\mid\mathcal G]$，并取 $[0,2\pi)$ 中一个可数稠密集 $\Theta$。对每个
$\theta\in\Theta$，实值情形给出
\[
 \Re(e^{-i\theta}Y)
 =\E[\Re(e^{-i\theta}X)\mid\mathcal G]
 \leq Z\quad\text{a.s.}
\]
删除这些例外集的可数并后，对同一个 $\omega$ 及全部 $\theta\in\Theta$ 上式成立。由于
$\sup_{\theta\in\Theta}\Re(e^{-i\theta}z)=|z|$，得到 $|Y|\leq Z$。

最后取期望并用定义中的 $G=\Omega$，得到
\[
 \|Y\|_1\leq\E Z=\E|X|.
\]
\end{proof}

\begin{proposition}[取出已知因子]\label{proposition:5-2-1-3}
设 $Z$ 为 $\mathcal G$-可测实值或复值随机变量。若 $Z$ 有界且 $X\in L^1$，则
\begin{equation}\label{eq:5-2-pull-out}
 \E[ZX\mid\mathcal G]=Z\E[X\mid\mathcal G].
\end{equation}
更一般地，只要 $X\in L^1$ 且 $ZX\in L^1$，同一等式仍成立，并且右端自动可积。
\end{proposition}

\begin{proof}
先设 $Z=\sum_{j=1}^mc_j\1_{A_j}$ 是有界 $\mathcal G$-简单函数，其中 $A_j$ 两两不交。
对任意 $G\in\mathcal G$，
\[
 \int_G Z\E[X\mid\mathcal G]\,\dd\Prob
 =\sum_{j=1}^mc_j\int_{G\cap A_j}\E[X\mid\mathcal G]\,\dd\Prob
 =\sum_{j=1}^mc_j\int_{G\cap A_j}X\,\dd\Prob
 =\int_GZX\,\dd\Prob.
\]
故唯一性给出 \eqref{eq:5-2-pull-out}。对有界实值 $Z$，取一致有界且逐点收敛于 $Z$ 的
$\mathcal G$-简单函数 $Z_n$。支配收敛定理分别用于 $Z_nX$ 与
$Z_n\E[X\mid\mathcal G]$，再结合 $L^1$ 收缩
\[
 \|\E[Z_nX-ZX\mid\mathcal G]\|_1\leq\|Z_nX-ZX\|_1,
\]
即可在 $L^1$ 中令 $n\to\infty$。复值 $Z$ 对实部、虚部分别处理。

现在只假设 $ZX\in L^1$。先设 $Z$ 实值，并令
\[
 Z_m=(-m)\vee(Z\wedge m).
\]
则 $|(Z_m-Z)X|\leq|ZX|\1_{\{|Z|>m\}}$，所以 $Z_mX\to ZX$ 于 $L^1$。由
\eqref{eq:5-2-l1-contraction}，
$\E[Z_mX\mid\mathcal G]\to\E[ZX\mid\mathcal G]$ 于 $L^1$。取一个几乎处处收敛的子列；
在该子列上，已证的有界因子公式给左边等于
$Z_m\E[X\mid\mathcal G]$，而后者逐点收敛到
$Z\E[X\mid\mathcal G]$。因此 \eqref{eq:5-2-pull-out} a.s. 成立，且右端因等于一个
$L^1$ 变量而可积。复值情形分解 $Z$ 的实部、虚部；由 $|\Re Z|,|\Im Z|\leq|Z|$，所需积
仍可积。
\end{proof}

\begin{theorem}[塔式性质]\label{theorem:5-2-1-4}
若 $\mathcal H\subset\mathcal G\subset\mathcal F$ 且 $X\in L^1$，则
\begin{equation}\label{eq:5-2-tower}
 \E\bigl[\E[X\mid\mathcal G]\mid\mathcal H\bigr]
 =\E[X\mid\mathcal H]\quad\text{a.s.}
\end{equation}
特别地，$\E\E[X\mid\mathcal G]=\E X$。
\end{theorem}

\begin{proof}
左端是 $\mathcal H$-可测可积变量。对任意 $H\in\mathcal H\subset\mathcal G$，两次使用
条件期望的积分刻画，得到
\[
 \int_H\E[\E[X\mid\mathcal G]\mid\mathcal H]\,\dd\Prob
 =\int_H\E[X\mid\mathcal G]\,\dd\Prob
 =\int_HX\,\dd\Prob.
\]
所以左端是 $\E[X\mid\mathcal H]$ 的版本，唯一性给出 \eqref{eq:5-2-tower}。令
$\mathcal H=\{\varnothing,\Omega\}$ 即得最后一式。
\end{proof}

\begin{example}[有限分割]\label{ex:5-2-1-1}
设 $A_1,\dots,A_m$ 是 $\Omega$ 的有限可测分割，$\mathcal G=\sigma(A_1,\dots,A_m)$。定义
\[
 Y=\sum_{i:\Prob(A_i)>0}
   \frac{\E[X\1_{A_i}]}{\Prob(A_i)}\1_{A_i},
\]
并在零概率块上任取有限常数。每个 $G\in\mathcal G$ 是若干 $A_i$ 的并；在这些块上逐项求和，
便有 $\int_GY=\int_GX$，所以 $Y=\E[X\mid\mathcal G]$。

回到\cref{ex:5-2-1-prediction-pressure}，奇数块上
\[
 \frac{\E[D^2\1_{A_o}]}{\Prob(A_o)}
 =\frac{(1+9+25)/6}{1/2}=\frac{35}{3},
\]
偶数块上同理为 $56/3$。若把被预测量换成 $D$，两块平均分别为
$(1+3+5)/3=3$ 与 $(2+4+6)/3=4$；“知道奇偶”并不会神奇地把平方赔付也变成 $3$ 与 $4$。
\end{example}

\begin{example}[独立信息]\label{ex:5-2-1-2}
若 $X\in L^1$ 且 $\sigma(X)$ 与 $\mathcal G$ 独立，则
\[
 \E[X\mid\mathcal G]=\E X\quad\text{a.s.}
\]
先对 $\sigma(X)$-可测简单函数 $X$，独立性给
$\E[X\1_G]=\E X\,\Prob(G)$，$G\in\mathcal G$。对一般实值或复值 $X\in L^1$，取
相应简单函数作 $L^1$ 逼近，并用
$|\E[(X-X_n)\1_G]|\leq\|X-X_n\|_1$ 令 $n\to\infty$。常数 $\E X$ 因而满足条件期望的
积分刻画。
\end{example}

\begin{example}[已知变量]\label{ex:5-2-1-3}
若 $X\in L^1$ 本身 $\mathcal G$-可测，则 $X$ 已满足定义中的两项条件，故
$\E[X\mid\mathcal G]=X$ a.s.。塔式性质因此可以理解为：先压缩到 $\mathcal G$，再压缩到更粗的
$\mathcal H$，与直接压缩到 $\mathcal H$ 得到同一结果。
\end{example}

\begin{figure}[htbp]
\centering
\begin{tikzpicture}[node distance=13mm and 18mm]
  \node[book strong node,minimum width=31mm] (x) {$X$\\$\mathcal F$-可测};
  \node[book strong node,right=of x,minimum width=38mm] (y)
    {$Y=\E[X\mid\mathcal G]$\\$\mathcal G$-可测};
  \node[book warm node,below=of y,minimum width=38mm] (tests)
    {$G\in\mathcal G$\\$\displaystyle\int_GY=\int_GX$};
  \draw[book accent arrow] (x)--node[above,book label,align=center,font=\scriptsize]{忘掉不可\\观察细节}(y);
  \draw[book arrow] (y)--(tests);
  \draw[book dashed arrow,bend right=28] (x.south)--node[below left,book label]{在全部可观察事件上保持积分}(tests.west);
\end{tikzpicture}
\caption{条件期望把随机变量压到较粗的信息层，但保留每个可观察事件上的积分。这里的箭头不是点态取值公式。}
\label{fig:5-2-1-information-compression}
\end{figure}

因子取出公式要求因子可观察；若 $Z$ 不是 $\mathcal G$-可测，右端甚至未必
$\mathcal G$-可测。一般无界因子还必须保证 $ZX\in L^1$。同样，版本只逐个 a.s. 唯一；未经
可数协调，不能把一族逐参数等式升级为一个同时成立的点态等式。

当信息层随时间递增时，塔式性质给出逐层压缩的一致性；取因子公式则把可观察量移出条件平均。
这两种机制共同保证逐步信息更新中条件平均计算的相容性。

\begin{exercise}[有限信息上的保费]
设 $D$ 为公平骰子，$X=\1_{\{D\geq4\}}$，$\mathcal G=\sigma(\{D\text{ 为奇数}\})$。
计算 $\E[X\mid\mathcal G]$，并在生成分割的每一块上核对积分。
\end{exercise}

\begin{exercise}[独立信息的截断证明]
从有界 $\sigma(X)$-简单函数开始，使用 $X\wedge n$ 及正负部分，完整证明
\cref{ex:5-2-1-2} 的实值情形。
\end{exercise}

\begin{exercise}[塔式性质与已知因子]
设 $\mathcal H\subset\mathcal G$，$Z$ 有界且 $\mathcal H$-可测。证明
\[
 \E[ZX\mid\mathcal H]=Z\E[X\mid\mathcal H],
 \qquad
 \E[Z\E[X\mid\mathcal G]]=\E[ZX].
\]
\end{exercise}

条件期望为每个测试函数给出一个条件平均。若状态空间是标准 Borel，这些平均可以在一个共同的
可数骨架上协调成随机概率测度；下一小节完成这项构造。


\subsection{正则条件分布、测度分解与 Markov 核}\label{subsec:5-2-2}
\index{正则条件分布}\index{测度分解}\index{后验核}

假设 $X$ 与 $Y$ 取离散值。给定 $Y=y$ 时，$X$ 的条件分布是一族概率
$A\mapsto\Prob(X\in A\mid Y=y)$；若 $Y$ 连续，这一条件概率族仍应存在，只是不能再由
$\Prob(Y=y)$ 作分母。我们需要的是一个关于 $y$ 可测的概率测度族。第~3~章已经把这样的族称为
概率核；现在的问题是，条件期望给出的逐事件等价类能否同时选成一个核。

下面两个计算先给出候选，并显示抽象结论必须恢复哪些熟悉公式。

\begin{example}[离散条件分布]\label{ex:5-2-2-1}
设 $Y$ 的值域包含于标准 Borel 空间 $T$ 的可数集 $C=\{y_1,y_2,\dots\}$，而
$X$ 取值于可测空间 $(S,\mathcal S)$。固定一个 $S$ 上的概率测度 $\mu_0$。定义
\begin{equation}\label{eq:5-2-discrete-kernel}
 k(y,A)=
 \begin{cases}
 \displaystyle\frac{\Prob(X\in A,Y=y)}{\Prob(Y=y)},
   &y\in C,\ \Prob(Y=y)>0,\\[7pt]
 \mu_0(A),&\text{其他情形}.
 \end{cases}
\end{equation}
标准 Borel 空间中的单点是 Borel 集，所以对每个 $A$，函数 $y\mapsto k(y,A)$ 是一个只在
可数集上偏离常数的 Borel 函数。固定 $y$ 时，$A\mapsto k(y,A)$ 是概率测度，故 $k$ 是概率核。
对任意 Borel 集 $B\subset T$，
\[
 \begin{aligned}
 \int_{\{Y\in B\}}k(Y,A)\,\dd\Prob
 &=\sum_{j:y_j\in B}\Prob(Y=y_j)k(y_j,A)\\
 &=\sum_{j:y_j\in B}\Prob(X\in A,Y=y_j)\\
 &=\Prob(X\in A,Y\in B).
 \end{aligned}
\]
零概率状态上的 $\mu_0$ 没有进入任何一项，说明这些状态的核值不由联合分布决定。
\end{example}

\begin{example}[联合密度]\label{ex:5-2-2-2}
设 $X\in\R^d$、$Y\in\R^m$ 的联合分布关于 Lebesgue 测度有密度 $p(x,y)$，并令
\[
 p_Y(y)=\int_{\R^d}p(x,y)\,\dd x.
\]
Tonelli 定理说明 $p_Y$ 可测。取任一概率密度 $r$ 于 $\R^d$，定义
\begin{equation}\label{eq:5-2-density-kernel}
 k(y,\dd x)=
 \begin{cases}
 \dfrac{p(x,y)}{p_Y(y)}\,\dd x,&p_Y(y)>0,\\[6pt]
 r(x)\,\dd x,&p_Y(y)=0.
 \end{cases}
\end{equation}
固定 Borel 集 $A$ 后，参数积分的可测性说明 $y\mapsto k(y,A)$ 可测。对 Borel 集
$B\subset\R^m$，Fubini--Tonelli 计算给出
\[
 \begin{aligned}
 \int_{\{Y\in B\}}k(Y,A)\,\dd\Prob
 &=\int_B k(y,A)p_Y(y)\,\dd y\\
 &=\int_B\int_Ap(x,y)\,\dd x\,\dd y
 =\Prob(X\in A,Y\in B).
 \end{aligned}
\]
因此密度比确实给出条件核；它是一般定义的一个后果，而不是一般定义本身。
\end{example}

\begin{remark}[标准 Borel 假设的边界]\label{ex:5-2-2-3}
若目标可测空间不是标准 Borel，正则条件分布可能不存在。本书不在这里构造所需的病态可测空间，
但下面的存在证明会逐处标出标准 Borel 假设的用途：它既提供可数生成的半直线骨架，又提供像为
Borel 集且逆映射可测的编码。只知道某个 $\sigma$-代数可数生成，或只知道一个可测函数族能够分离
点，并不能自动提供这样的双向编码。条件信息 $\mathcal G$ 是否可数生成则不是本定理的假设。
\end{remark}

\begin{example}[确定性条件核]\label{ex:5-2-2-4}
若 $X=h(Y)$，其中 $h:T\to S$ 可测，则
\[
 k(y,\cdot)=\delta_{h(y)}
\]
是从 $T$ 到 $S$ 的概率核。对 $A\in\mathcal S$，
$k(Y,A)=\1_A(h(Y))=\1_{\{X\in A\}}$，故它已经是相应指标函数关于 $Y$ 的条件期望；并且
对非负 Borel 函数 $f$，
\[
 \int_Sf(x)k(y,\dd x)=f(h(y)).
\]
这个公式在 $\Prob(Y=y)=0$ 时仍有确定含义，不过在 $Y$ 的分布零集上改动整个核也不会改变任何
条件积分。
\end{example}

\begin{definition}[正则条件分布]\label{def:5-2-2-1}
设 $X:\Omega\to(S,\mathcal S)$ 是随机元，$\mathcal G\subset\mathcal F$。若从
$(\Omega,\mathcal G)$ 到 $(S,\mathcal S)$ 的概率核 $K$ 满足：对每个 $A\in\mathcal S$，
\begin{equation}\label{eq:5-2-rcd-def}
 K(\cdot,A)=\E[\1_{\{X\in A\}}\mid\mathcal G]
 \quad\text{a.s.},
\end{equation}
则称 $K$ 是 $X$ 给定 $\mathcal G$ 的一个\emph{正则条件分布}，记作
$K=\Law(X\mid\mathcal G)$。
\end{definition}

\begin{definition}[给定随机变量的核]\label{def:5-2-2-2}
设 $S,T$ 都是标准 Borel 空间，$X:\Omega\to S$、$Y:\Omega\to T$。若从 $T$ 到 $S$ 的
概率核 $k$ 满足
\[
 k(Y(\omega),A)
 =\E[\1_{\{X\in A\}}\mid\sigma(Y)](\omega)
 \quad\text{a.s.}
\]
对每个 $A$ 成立，则把 $k(y,\cdot)$ 记作 $\Law(X\mid Y=y)$。这个点值记号包含了版本选择；
它不表示对零概率事件使用比值。
\end{definition}

\begin{definition}[测度分解]\label{def:5-2-2-3}
设 $\pi$ 是 $(S\times T,\mathcal S\otimes\mathcal T)$ 上的概率，第一边缘为 $\mu$。若从
$S$ 到 $T$ 的概率核 $K$ 满足
\begin{equation}\label{eq:5-2-disintegration-rectangles}
 \pi(A\times B)=\int_AK(x,B)\,\mu(\dd x)
 \qquad(A\in\mathcal S,\ B\in\mathcal T),
\end{equation}
就把它写作
\[
 \pi(\dd x,\dd y)=\mu(\dd x)K(x,\dd y),
\]
并称这是 $\pi$ 相对于第一边缘的\emph{测度分解}（disintegration）。
\end{definition}

逐个 $A$ 选择 \eqref{eq:5-2-rcd-def} 的版本不够：$\mathcal S$ 通常不可数，而每个版本各有一个
例外零集。若只在一个可数集合代数上选版本，再声称逐点作 \Caratheodory{} 扩张，也仍有缺口，
因为这个代数拥有不可数多条下降序列，不能用一个可数零集保证逐点可数可加。实线上的分布函数把
可数协调压缩为有理阈值间的单调性与右连续关系，正好绕开这个障碍。

\begin{theorem}[正则条件分布的存在性]\label{theorem:5-2-2-1}
设 $(S,\mathcal S)$ 是标准 Borel 空间，$X:\Omega\to S$。对任意子 $\sigma$-代数
$\mathcal G\subset\mathcal F$，都存在 $\Law(X\mid\mathcal G)$ 的正则版本；这里不要求
$\mathcal G$ 可数生成。

若 $Y:\Omega\to T$ 还取值于标准 Borel 空间，则存在从 $T$ 到 $S$ 的概率核 $k$，使
$k(Y,\cdot)$ 是 $\Law(X\mid\sigma(Y))$ 的正则版本。
\end{theorem}

\begin{proof}
\emph{第一步：把目标编码进实线。}
由标准 Borel 编码定理~\ref{theorem:1-4-2-4}，存在 Borel 集 $B\subset[0,1]$ 与 Borel
同构 $\iota:S\to B$。令 $Z=\iota(X)$。对每个 $q\in\Q$，选择一个版本
\[
 f_q=\E[\1_{\{Z\leq q\}}\mid\mathcal G],
\]
并把它截到 $[0,1]$；截断只改变零测集。

若 $q<r$，条件期望的单调性给 $f_q\leq f_r$ a.s.。若 $q<0$，则 $f_q=0$ a.s.；若
$q\geq1$，则 $f_q=1$ a.s.。还需协调右连续性。对固定 $q$，取有理数
$r_{q,n}\downarrow q$。删去单调性例外集后，$f_{r_{q,n}}$ 逐点非增；记其极限为 $h_q$。对
$G\in\mathcal G$，支配收敛定理和测度从上连续性给出
\[
 \begin{aligned}
 \int_Gh_q\,\dd\Prob
 &=\lim_n\int_Gf_{r_{q,n}}\,\dd\Prob
 =\lim_n\Prob(G\cap\{Z\leq r_{q,n}\})\\
 &=\Prob(G\cap\{Z\leq q\}).
 \end{aligned}
\]
所以 $h_q$ 与 $f_q$ 是同一个条件期望版本。把 $\{f_q\ne\lim_n f_{r_{q,n}}\}$、全部有理对
$q<r$ 的单调性失败集
以及端点关系的失败集取可数并，得到一个 $N\in\mathcal G$、$\Prob(N)=0$。在
$N$ 上把
$f_q$ 全部改成 $\1_{\{q\geq0\}}$；于是对每个 $\omega$ 都有
\begin{equation}\label{eq:5-2-rational-cdf-relations}
 0\leq f_q(\omega)\leq f_r(\omega)\leq1\quad(q<r),
 \qquad
 f_q(\omega)=\inf_{r>q,\ r\in\Q}f_r(\omega),
\end{equation}
并且 $f_q=0$ 于 $q<0$、$f_q=1$ 于 $q\geq1$。
最后一个等式使用了 $(f_r)_{r\in\Q}$ 的单调性：对任意有理数 $r>q$，最终都有
$r_{q,n}<r$，故沿 $(r_{q,n})$ 的极限正是对全部有理数 $r>q$ 所取的下确界。

\emph{第二步：逐点生成概率测度。}
对 $t\in\R$ 定义
\begin{equation}\label{eq:5-2-cdf-extension}
 F_\omega(t)=\inf_{q>t,\ q\in\Q}f_q(\omega).
\end{equation}
它非减。若 $t_n\downarrow t$，则 $F_\omega(t_n)\geq F_\omega(t)$；另一方面，对任意有理
$q>t$，充分大的 $n$ 满足 $t_n<q$，因而
$F_\omega(t_n)\leq f_q(\omega)$。先令 $n\to\infty$，再对 $q>t$ 取下确界，得到
$F_\omega(t_n)\downarrow F_\omega(t)$。所以 $F_\omega$ 右连续。端点关系还给出
\[
 F_\omega(t)=0\quad(t<0),
 \qquad F_\omega(t)=1\quad(t\geq1),
\]
且由 \eqref{eq:5-2-rational-cdf-relations}，$F_\omega(q)=f_q(\omega)$ 对每个有理 $q$ 成立。
Lebesgue--Stieltjes 构造~\ref{theorem:2-3-3-2} 因而给出实线上的概率测度
$L(\omega,\cdot)$，其分布函数为 $F_\omega$，并且每个 $L(\omega,\cdot)$ 集中在 $[0,1]$。

\emph{第三步：验证参数可测性与条件恒等式。}
对每个 $t$，式 \eqref{eq:5-2-cdf-extension} 是可数个 $\mathcal G$-可测函数的下确界，故
$\omega\mapsto L(\omega,(-\infty,t])$ 可测。令
\[
 \mathcal C=\{A\in\Borel(\R):\omega\mapsto L(\omega,A)
 \text{ 是 }\mathcal G\text{-可测的}\}.
\]
$\mathcal C$ 含实线，对补集封闭，并对两两不交的可数并封闭；后一项由逐点测度可加性和可测函数
部分和的极限得到。因此 $\mathcal C$ 是 $\lambda$-系统。它包含所有有理端点左半直线，这些集合
构成生成 $\Borel(\R)$ 的 $\pi$-系统。$\pi$--$\lambda$ 定理给
$\mathcal C=\Borel(\R)$，故 $L$ 是概率核。

再令
\[
 \mathcal D=\{A\in\Borel(\R):L(\cdot,A)
 \text{ 是 }\E[\1_{\{Z\in A\}}\mid\mathcal G]\text{ 的版本}\}.
\]
由 $L(\cdot,(-\infty,q])=f_q$，所有有理端点左半直线属于 $\mathcal D$。常数一处理全空间；
线性处理补集；若 $A_n$ 两两不交且属于 $\mathcal D$，则对 $G\in\mathcal G$，单调收敛给
\[
 \int_G L\!\left(\cdot,\bigcup_nA_n\right)\,\dd\Prob
 =\sum_n\int_GL(\cdot,A_n)\,\dd\Prob
 =\sum_n\Prob(G\cap\{Z\in A_n\}).
\]
所以 $\mathcal D$ 也是 $\lambda$-系统。再次使用 $\pi$--$\lambda$ 定理，得到
$\mathcal D=\Borel(\R)$。

特别地，$L(\cdot,B)$ 是 $\E[\1_{\{Z\in B\}}\mid\mathcal G]=1$ 的版本。故
$N_B=\{L(\cdot,B)\ne1\}\in\mathcal G$ 是零测集。固定 $b_0\in B$，在 $N_B$ 上把
$L(\omega,\cdot)$ 改成 $\delta_{b_0}$。修改后的核在每个 $\omega$ 都集中于 $B$，并且仍满足
全部条件恒等式。最后定义
\[
 K(\omega,A)=L(\omega,\iota(A)),\qquad A\in\mathcal S.
\]
因为 $\iota$ 及其逆映射都可测，$\iota(A)$ 是 $B$ 中、从而是 $\R$ 中的 Borel 集；所以 $K$
是从 $(\Omega,\mathcal G)$ 到 $S$ 的概率核，并满足 \eqref{eq:5-2-rcd-def}。

\emph{第四步：给定 $Y=y$ 的版本。}
现在令 $\mathcal G=\sigma(Y)$。在第一步只出现可数个函数 $f_q$。对每个 $q$ 使用
Doob--Dynkin 因子化，先写成 $f_q=h_q(Y)$；因 $f_q\in[0,1]$ a.s.，再把 $h_q$ 截断到
$[0,1]$，仍保持该等式 a.s. 成立。于是可设 $h_q:T\to[0,1]$ Borel 可测。记
$\nu=\Law(Y)$，并固定 $b_0\in B$。由 $f_q$ 已经满足
\eqref{eq:5-2-rational-cdf-relations}，下列 Borel 集的 $\nu$-测度都为零：
\begin{align*}
 &\{y:h_q(y)>h_r(y)\}\quad(q<r),\\
 &\left\{y:h_q(y)\ne\inf_{r>q,\ r\in\Q}h_r(y)\right\}\quad(q\in\Q),\\
 &\{y:h_q(y)\ne0\}\quad(q<0),
 \qquad
 \{y:h_q(y)\ne1\}\quad(q\ge1).
\end{align*}
把它们的可数并记为 $E$。在 $E$ 上同时改定
\[
 h_q(y)=\1_{\{b_0\le q\}},
\]
则 $h_q(Y)=f_q$ 仍对所有有理数 $q$ 同时 a.s. 成立，且上述单调、端点和右连续关系
对每个 $y\in T$ 成立。

对 $y\in T$、$t\in\R$ 定义
\[
 F_y(t)=\inf_{q>t,\ q\in\Q}h_q(y).
\]
第二步的逐式论证说明 $F_y$ 非降、右连续，且
$F_y(-\infty)=0$、$F_y(+\infty)=1$。因此它唯一决定概率测度
$\ell(y,\cdot)$，并有
\[
 \ell(y,(-\infty,q])=h_q(y)\qquad(q\in\Q).
\]
对每个实数 $t$，$y\mapsto F_y(t)$ 是可数个 Borel 函数的下确界。把使
$y\mapsto\ell(y,A)$ 可测的 Borel 集 $A\subset\R$ 组成 $\lambda$-系统，并从有理端点左半直线出发使用
$\pi$--$\lambda$ 定理，可得 $\ell$ 是从 $T$ 到 $\R$ 的概率核。另一次同样的
$\lambda$-系统计算由等式
\[
 \ell(Y,(-\infty,q])=f_q\quad\text{a.s.}
\]
出发，得到对每个 $A\in\Borel(\R)$，
\[
 \ell(Y,A)=\E[\1_{\{Z\in A\}}\mid\sigma(Y)]
 \quad\text{a.s.}
\]
特别地，$\ell(Y,B)=1$ a.s.，所以
$E_B=\{y:\ell(y,B)\ne1\}$ 满足 $\nu(E_B)=0$。在 $E_B$ 上把整个测度
$\ell(y,\cdot)$ 改为 $\delta_{b_0}$，得到的核对每个 $y$ 都集中于 $B$，且不改变任何一条 a.s. 条件恒等式。
最后定义
\[
 k(y,A)=\ell(y,\iota(A)),\qquad A\in\mathcal S.
\]
则 $k$ 是所需的从 $T$ 到 $S$ 的概率核。这里因子化的是有理阈值的可数族，而不是尚未构造的“整个核”。
\end{proof}

\begin{remark}[正则条件分布构造的三个步骤]\label{rem:5-2-2-existence-checkpoints}
上述构造的三个步骤各有独立作用：有理阈值把版本选择压成可数问题；逐点分布函数保证每个参数处得到的
确是可数可加概率；Borel 同构的可测逆把实线上的核安全地送回原状态空间。

在\cref{ex:5-2-2-1}中，有理阈值条件期望就是每个正概率分割块上的离散累计概率，生成的测度恢复
\eqref{eq:5-2-discrete-kernel}。在\cref{ex:5-2-2-2}中，核在每个左半直线上的值由
\eqref{eq:5-2-density-kernel} 积分得到；测度唯一性因而恢复密度比。标准 Borel 假设作用在目标
空间 $S$，不是条件信息 $\mathcal G$。
\end{remark}

一旦核已经存在，所有 Borel 测试函数可以由指标函数和单调收敛统一处理。

\begin{theorem}[核表示条件期望]\label{theorem:5-2-2-2}
若 $K=\Law(X\mid\mathcal G)$，则对每个非负 Borel 函数 $f:S\to[0,\infty]$，
\begin{equation}\label{eq:5-2-kernel-ce}
 \E[f(X)\mid\mathcal G](\omega)
 =\int_Sf(x)K(\omega,\dd x)\quad\text{a.s.}
\end{equation}
两边允许为 $+\infty$。若 $f$ 为实值或复值且 $f(X)\in L^1$，同一公式仍成立；对 a.e.
$\omega$，函数 $f$ 关于 $K(\omega,\cdot)$ 绝对可积。
\end{theorem}

\begin{proof}
对 $f=\1_A$，公式就是正则条件分布的定义。非负简单函数情形由有限线性性得到。对一般
$f\geq0$，取非负简单函数 $s_n\uparrow f$。固定 $\omega$ 后，对 $K(\omega,\cdot)$ 使用单调
收敛；同时对扩展条件期望使用 \eqref{eq:5-2-conditional-mct}，得到
\[
 \int s_n\,\dd K(\omega,\cdot)\uparrow\int f\,\dd K(\omega,\cdot),
 \qquad
 \E[s_n(X)\mid\mathcal G]\uparrow\E[f(X)\mid\mathcal G]
\]
于同一个可数例外集之外，故 \eqref{eq:5-2-kernel-ce} 成立。

若 $f(X)\in L^1$，先把非负结论用于 $|f|$ 并取期望：
\[
 \E\int_S|f(x)|K(\cdot,\dd x)=\E|f(X)|<\infty.
\]
所以核积分对 a.e. $\omega$ 绝对收敛。实值函数分成正负部，复值函数再分实部、虚部，利用条件
期望线性即可。
\end{proof}

\begin{theorem}[测度分解定理]\label{theorem:5-2-2-3}
设 $S,T$ 是标准 Borel 空间，$\pi$ 是 $S\times T$ 上的概率，第一边缘为 $\pi_S$。则存在从
$S$ 到 $T$ 的概率核 $K$，使
\[
 \pi(\dd x,\dd y)=\pi_S(\dd x)K(x,\dd y).
\]
若 $K'$ 也是这样的核，则
$K(x,\cdot)=K'(x,\cdot)$ 对 $\pi_S$-a.e. $x$ 成立。
\end{theorem}

\begin{proof}
在规范概率空间 $(S\times T,\mathcal S\otimes\mathcal T,\pi)$ 上令
$U(x,y)=x$、$V(x,y)=y$。由\cref{theorem:5-2-2-1}及其中的因子化部分，存在从 $S$ 到
$T$ 的概率核 $K$，使
$K(U,B)=\E[\1_{\{V\in B\}}\mid\sigma(U)]$ a.s.。对 $A\in\mathcal S$、
$B\in\mathcal T$，已知因子公式与塔式性质给出
\[
 \begin{aligned}
 \pi(A\times B)
 &=\E_\pi[\1_A(U)\1_B(V)]\\
 &=\E_\pi[\1_A(U)K(U,B)]
 =\int_AK(x,B)\,\pi_S(\dd x),
 \end{aligned}
\]
故矩形公式成立。

为证唯一性，由标准 Borel 编码取 $T$ 的一个可数生成 $\pi$-系统 $\mathcal P$。若 $K,K'$ 都满足
矩形公式，则对每个 $B\in\mathcal P$ 与每个 $A\in\mathcal S$，
\[
 \int_AK(x,B)\,\pi_S(\dd x)=\int_AK'(x,B)\,\pi_S(\dd x).
\]
Radon--Nikodym 密度的唯一性说明 $K(\cdot,B)=K'(\cdot,B)$ a.e.。对可数个
$B\in\mathcal P$ 合并零集；在其补集上，两个概率测度在生成 $\pi$-系统上相等，有限测度唯一性
定理~\ref{theorem:2-2-1-2} 便给出它们在 $\mathcal T$ 上处处相等。
\end{proof}

给定一个边缘和一个概率核，反向确实可以构造联合测度，而不仅是形式地写两个积分号。

\begin{proposition}[核重构联合分布]\label{proposition:5-2-2-4}
设 $X:\Omega\to(S,\mathcal S)$、$Y:\Omega\to(T,\mathcal T)$ 是随机元，
$\mu=\Law(X)$，并设从 $S$ 到 $T$ 的概率核 $K$ 满足：对每个
$B\in\mathcal T$，
\[
 K(X,B)=\E[\1_{\{Y\in B\}}\mid\sigma(X)]\quad\text{a.s.}
\]
则
$(X,Y)$ 的联合分布为 $\mu K$，即
\begin{equation}\label{eq:5-2-joint-from-kernel}
 \Prob(X\in A,Y\in B)=\int_AK(x,B)\,\mu(\dd x)
 \qquad(A\in\mathcal S,\ B\in\mathcal T).
\end{equation}
对每个非负或关于联合分布可积的可测函数 $f$，
\begin{equation}\label{eq:5-2-kernel-fubini}
 \E f(X,Y)=\int_S\int_T f(x,y)K(x,\dd y)\mu(\dd x).
\end{equation}
\end{proposition}

\begin{proof}
先说明 $\mu K$ 的含义。对 $E\in\mathcal S\otimes\mathcal T$，令 $E_x=\{y:(x,y)\in E\}$。
满足 $x\mapsto K(x,E_x)$ 可测的集合 $E$ 构成一个 $\lambda$-系统：全空间对应常数一，补集对应
$1-K(x,E_x)$，两两不交可数并对应非负可测函数部分和的极限。它包含全部矩形，因为
\[
 K(x,(A\times B)_x)=\1_A(x)K(x,B).
\]
矩形是生成积 $\sigma$-代数的 $\pi$-系统，故 $x\mapsto K(x,E_x)$ 对每个积可测 $E$ 都可测。
定义
\[
 (\mu K)(E)=\int_SK(x,E_x)\,\mu(\dd x).
\]
若 $E_n$ 两两不交，则其每个 $x$-截面也两两不交；核的可数可加性和单调收敛定理说明
$\mu K$ 可数可加。又 $(\mu K)(S\times T)=1$，故这是联合概率。指标函数、非负简单函数和单调
收敛依次给出
\[
 \int f\,\dd(\mu K)=\int_S\int_Tf(x,y)K(x,\dd y)\mu(\dd x)
\]
对全部 $f\geq0$ 成立。

现在利用条件分布。对矩形 $A\times B$，
\[
 \begin{aligned}
 \Prob(X\in A,Y\in B)
 &=\E\bigl[\1_A(X)\E[\1_B(Y)\mid X]\bigr]\\
 &=\E[\1_A(X)K(X,B)]
 =\int_AK(x,B)\,\mu(\dd x).
 \end{aligned}
\]
两个概率测度在生成矩形 $\pi$-系统上相等，故联合分布等于 $\mu K$。非负
\eqref{eq:5-2-kernel-fubini} 已由上段得到；有符号或复值情形先把公式用于 $|f|$，确认两次积分
绝对有限，再分正负部或实虚部。
\end{proof}

\begin{proposition}[核版本的共同零集唯一性]\label{proposition:5-2-2-5}
设 $S$ 为标准 Borel 空间。若 $K,K'$ 都是 $\Law(X\mid\mathcal G)$ 的正则版本，则存在
$N\in\mathcal G$、$\Prob(N)=0$，使对每个 $\omega\notin N$，
\[
 K(\omega,\cdot)=K'(\omega,\cdot)
\]
作为 $S$ 上的概率测度成立。
\end{proposition}

\begin{proof}
标准 Borel 编码给出 $\mathcal S$ 的一个可数生成 $\pi$-系统 $\mathcal P$。对每个
$A\in\mathcal P$，$K(\cdot,A)$ 与 $K'(\cdot,A)$ 都是同一个条件期望的版本，故二者在某个零集
外相等。合并这些可数个零集得到 $N\in\mathcal G$、$\Prob(N)=0$。固定
$\omega\notin N$，两个概率测度 $K(\omega,\cdot)$、$K'(\omega,\cdot)$ 在
$\mathcal P$ 上相等；有限测度唯一性把等式推广到全部 $A\in\mathcal S$。

这个共同零集不能删除。\cref{ex:5-2-2-4}中，在 $Y$ 的分布零集上把
$\delta_{h(y)}$ 换成任一概率测度，仍得到同一个条件分布版本。
\end{proof}

\begin{figure}[htbp]
\centering
\begin{tikzpicture}[node distance=13mm and 18mm]
  \node[book strong node,minimum width=27mm] (mu) {第一边缘\\$\mu(\dd x)$};
  \node[book strong node,right=of mu,minimum width=34mm] (cloud)
    {给定 $x$ 的条件概率\\$K(x,\dd y)$};
  \node[book warm node,right=of cloud,minimum width=34mm] (joint)
    {联合概率\\$\mu(\dd x)K(x,\dd y)$};
  \draw[book accent arrow] (mu)--node[above,book label]{指定条件核}(cloud);
  \draw[book accent arrow] (cloud)--node[above,book label]{对 $x$ 积分}(joint);
  \foreach \p in {-.34,0,.34}{
    \fill[BookAccent!72] ($(cloud.center)+(0.72,\p)$) circle (1.25pt);
  }
  \draw[book dashed arrow,bend right=28] (joint.south)--node[below,book label]{取第一边缘}(mu.south);
\end{tikzpicture}
\caption{测度分解先保留第一边缘，再在每个 $x$ 上附加条件概率；把这些概率测度积分起来便重构联合分布。}
\label{fig:5-2-2-disintegration-clouds}
\end{figure}

正则条件分布只在边缘 a.e. 意义下唯一，符号 $\Law(X\mid Y=y)$ 因而必须连同版本理解。
本小节的存在定理依赖目标状态空间标准 Borel；它没有断言任意可测空间都存在条件核。连续时间路径
分布的逐步构造还需要后面的扩张定理，此处不使用也不预设 Ionescu--Tulcea 定理。

\begin{exercise}[离散核的零概率行]
在\cref{ex:5-2-2-1}中把零概率状态上的 $\mu_0$ 换成另一概率 $\mu_1$。证明两个核给出同一个
关于 $Y$ 的条件分布版本，并找出它们可能不同的全部状态。
\end{exercise}

\begin{exercise}[密度模型的测度分解]
从 \eqref{eq:5-2-density-kernel} 出发，证明
\[
 p(x,y)\,\dd x\,\dd y
 =p_Y(y)\,\dd y\,k(y,\dd x),
\]
并把公式推广到任意非负 Borel 测试函数。
\end{exercise}

\begin{exercise}[核的共同零集]
完整写出\cref{proposition:5-2-2-5}中可数生成 $\pi$-系统：先把 $S$ Borel 编码为
$B\subset[0,1]$，再使用有理端点左半直线与 $B$ 的交。
\end{exercise}

条件核给出每个状态下的完整分布，条件期望则选取其中的积分。下一小节把这种积分预测同凸性、
平方误差几何和贝叶斯换测度公式联系起来。


\subsection{条件 Jensen、\texorpdfstring{$L^2$}{L2} 投影与贝叶斯公式}\label{subsec:5-2-3}
\index{条件 Jensen 不等式}\index{最小均方预测}\index{贝叶斯公式}

在\cref{ex:5-2-1-prediction-pressure}中，两块上的条件平均恰好最小化平方误差。这个现象不是有限
分割的巧合。条件期望一方面受凸性控制，另一方面在 $L^2$ 中与所有可观察方向正交。贝叶斯公式则
回答另一个表面上不同的问题：更换概率测度以后，同一信息层上的条件平均怎样改变？三者最终都只
使用条件期望的积分刻画。

\begin{definition}[信息子空间]\label{def:5-2-3-1}
定义
\[
 L^2(\mathcal G)=\{Y\in L^2(\Prob):Y\text{ 有一个 }\mathcal G\text{-可测版本}\}.
\]
\end{definition}

这是 $L^2(\mathcal F)$ 的闭线性子空间。线性性直接成立。为证闭性，设
$Y_n\in L^2(\mathcal G)$ 且
$Y_n\to Y$ 于 $L^2$，为每个 $Y_n$ 选取一个 $\mathcal G$-可测版本，再取子列使
\[
 \|Y_{n_k}-Y\|_2^2\leq 2^{-3k}\qquad(k\geq1).
\]
Markov 不等式给出
\[
 \sum_{k=1}^\infty
 \Prob\bigl(|Y_{n_k}-Y|>2^{-k}\bigr)
 \leq\sum_{k=1}^\infty2^{-k}<\infty.
\]
由 Borel--Cantelli I，$Y_{n_k}\to Y$ a.s.。令
\[
 C=\{\omega:(Y_{n_k}(\omega))_{k\geq1}\text{ 在 }\C\text{ 中为 Cauchy 列}\}.
\]
$C\in\mathcal G$，因为 Cauchy 条件可写成关于 $|Y_{n_j}-Y_{n_k}|$ 的可数次交并；函数
\[
 \widetilde Y(\omega)=
 \begin{cases}
  \lim_{k\to\infty}Y_{n_k}(\omega),&\omega\in C,\\
  0,&\omega\notin C
 \end{cases}
\]
是 $\mathcal G$-可测，并且 $\widetilde Y=Y$ a.s.，故 $Y\in L^2(\mathcal G)$。这个论证同时
适用于实值与复值 $L^2$。

要把 $\E[X\mid\mathcal G]$ 解释为这个子空间上的投影，还须验证它保持平方可积。
\cref{theorem:5-2-3-1} 中取 $\varphi(u)=u^2$ 将给出所需的 $L^2$ 收缩估计。

\begin{definition}[后验核]\label{def:5-2-3-2}
设参数空间 $\Theta$ 与观测空间 $T$ 为标准 Borel 空间，$\pi(\dd\theta)$ 是先验概率，
$L(\theta,\dd y)$ 是似然核。联合分布
\[
 \pi(\dd\theta)L(\theta,\dd y)
\]
相对于观测边缘的测度分解核 $\pi(\dd\theta\mid y)$ 称为\emph{后验核}。它在观测边缘
a.e. 意义下唯一；零边缘质量的观测值上仍需选择版本。
\end{definition}

\begin{definition}[换测度]\label{def:5-2-3-3}
设 $\mathbb Q$ 是 $(\Omega,\mathcal F)$ 上的概率且 $\mathbb Q\ll\Prob$，记
\[
 Z=\frac{\dd\mathbb Q}{\dd\Prob}.
\]
称这是从 $\Prob$ 到 $\mathbb Q$ 的密度或似然比。对信息层 $\mathcal G$，记
\[
 Z_{\mathcal G}=\E_\Prob[Z\mid\mathcal G].
\]
\end{definition}

\begin{proposition}[限制信息上的密度]
在定义~\ref{def:5-2-3-3} 的条件下，$Z_{\mathcal G}$ 是
$\mathbb Q|_{\mathcal G}$ 关于 $\Prob|_{\mathcal G}$ 的 Radon--Nikodym 导数。
\end{proposition}

\begin{proof}
对每个 $G\in\mathcal G$，
\[
 \mathbb Q(G)=\int_GZ\,\dd\Prob=\int_GZ_{\mathcal G}\,\dd\Prob,
\]
而 $Z_{\mathcal G}$ 非负且 $\mathcal G$-可测。这正是所述 Radon--Nikodym 导数的定义。
\end{proof}

先建立凸性估计。证明中只用一维凸函数的支撑直线，并把不可数条直线缩成可数族。

\begin{theorem}[条件 Jensen]\label{theorem:5-2-3-1}
设 $\varphi:\R\to\R$ 为有限凸函数。若实值 $X$ 满足
$X,\varphi(X)\in L^1$，则
\begin{equation}\label{eq:5-2-conditional-jensen}
 \varphi\bigl(\E[X\mid\mathcal G]\bigr)
 \leq\E[\varphi(X)\mid\mathcal G]
 \quad\text{a.s.}
\end{equation}
左端也属于 $L^1$。
\end{theorem}

\begin{proof}
有限凸函数在实线上连续。对每个 $r\in\Q$，凸割线斜率的单调性给出一个实数 $a_r$，使
\begin{equation}\label{eq:5-2-convex-support}
 \ell_r(u):=\varphi(r)+a_r(u-r)\leq\varphi(u)
 \qquad(u\in\R).
\end{equation}
例如可在左、右导数形成的闭区间
$[\varphi'_-(r),\varphi'_+(r)]$ 中选取 $a_r$；
\eqref{eq:5-2-convex-support} 直接来自左右割线斜率的次序。

这个可数族恢复 $\varphi$：
\begin{equation}\label{eq:5-2-countable-affine-envelope}
 \varphi(u)=\sup_{r\in\Q}\ell_r(u).
\end{equation}
为证反向不等式，固定 $u$ 并取有理数 $r_n\to u$。当 $r_n\in(u-1,u+1)$ 时，凸割线斜率的
单调性给出明确的界
\[
 \varphi(u-1)-\varphi(u-2)
 \leq a_{r_n}
 \leq \varphi(u+2)-\varphi(u+1).
\]
故 $(a_{r_n})$ 有界。由 $\varphi(r_n)\to\varphi(u)$，
\[
 \ell_{r_n}(u)=\varphi(r_n)+a_{r_n}(u-r_n)\longrightarrow\varphi(u).
\]
结合每条支撑线均不超过 $\varphi$，便得 \eqref{eq:5-2-countable-affine-envelope}。

令 $Y=\E[X\mid\mathcal G]$。对每个 $r\in\Q$，线性与正性给出
\[
 \ell_r(Y)
 =\E[\ell_r(X)\mid\mathcal G]
 \leq\E[\varphi(X)\mid\mathcal G]
 \quad\text{a.s.}
\]
合并可数个例外集，再对 $r$ 取上确界并使用
\eqref{eq:5-2-countable-affine-envelope}，得到
\eqref{eq:5-2-conditional-jensen}。最后固定任一支撑线 $\ell_{r_0}$。我们已经得到
\[
 \ell_{r_0}(Y)\leq\varphi(Y)\leq\E[\varphi(X)\mid\mathcal G].
\]
左右两端都可积。更明确地，
$\varphi(Y)^+\leq\E[\varphi(X)\mid\mathcal G]^+$ 且
$\varphi(Y)^-\leq\ell_{r_0}(Y)^-$，所以 $\varphi(Y)\in L^1$。
\end{proof}

\begin{remark}[从凸性进入平方几何]\label{rem:5-2-3-l2-entry}
若实值 $X\in L^2$，在\cref{theorem:5-2-3-1}中取 $\varphi(u)=u^2$，得到
\begin{equation}\label{eq:5-2-l2-contraction}
 \bigl|\E[X\mid\mathcal G]\bigr|^2
 \leq\E[|X|^2\mid\mathcal G],
 \qquad
 \|\E[X\mid\mathcal G]\|_2\leq\|X\|_2.
\end{equation}
所以 $P_{\mathcal G}X:=\E[X\mid\mathcal G]$ 确实属于
$L^2(\mathcal G)$。对\cref{ex:5-2-1-prediction-pressure}，
$P_{\mathcal G}(D^2)$ 正是奇数块上的 $35/3$ 与偶数块上的 $56/3$。残差同每个可观察的
$L^2$ 方向正交后，平方误差的最小化性质将直接由勾股恒等式得到。
\end{remark}

复值 $X\in L^2$ 也满足 \eqref{eq:5-2-l2-contraction}：由
$|\E[X\mid\mathcal G]|\leq\E[|X|\mid\mathcal G]$，再把实值条件 Jensen 用于非负变量
$|X|$ 与函数 $u\mapsto u^2$ 即得。

\begin{proposition}[条件 \Holder{}]\label{proposition:5-2-3-holder}
设 $1<p,q<\infty$、$p^{-1}+q^{-1}=1$，$U\in L^p$、$V\in L^q$。则
\begin{equation}\label{eq:5-2-conditional-holder}
 \bigl|\E[UV\mid\mathcal G]\bigr|
 \leq
 \E[|U|^p\mid\mathcal G]^{1/p}
 \E[|V|^q\mid\mathcal G]^{1/q}
 \quad\text{a.s.}
\end{equation}
\end{proposition}

\begin{proof}
普通 \Holder{} 不等式~\ref{theorem:3-3-2-2} 先给 $UV\in L^1$。记
\[
 A=\E[|U|^p\mid\mathcal G],
 \qquad B=\E[|V|^q\mid\mathcal G].
\]
$A,B$ 非负且有限 a.s.。固定 $\varepsilon>0$，对点态 Young 不等式应用随机归一化
\[
 a=\frac{|U|}{(A+\varepsilon)^{1/p}},
 \qquad
 b=\frac{|V|}{(B+\varepsilon)^{1/q}},
\]
得到
\[
 \frac{|UV|}{(A+\varepsilon)^{1/p}(B+\varepsilon)^{1/q}}
 \leq
 \frac{|U|^p}{p(A+\varepsilon)}
 +\frac{|V|^q}{q(B+\varepsilon)}.
\]
所有分母倒数都是有界的 $\mathcal G$-可测变量，可以从条件期望中取出。于是
\[
 \frac{\E[|UV|\mid\mathcal G]}
 {(A+\varepsilon)^{1/p}(B+\varepsilon)^{1/q}}
 \leq\frac{A}{p(A+\varepsilon)}+\frac{B}{q(B+\varepsilon)}\leq1.
\]
依次取 $\varepsilon=1/n$，合并这可数个不等式的例外集，再令 $n\to\infty$。零分母处也由
该极限给出正确的零值，得到
\[
 \E[|UV|\mid\mathcal G]\leq A^{1/p}B^{1/q}.
\]
再用
$|\E[UV\mid\mathcal G]|\leq\E[|UV|\mid\mathcal G]$ 即得
\eqref{eq:5-2-conditional-holder}。整个论证没有在 $A=0$ 或 $B=0$ 处作除法。
\end{proof}

\begin{example}[在有限信息块上核对条件 \Holder{}]
令 $D$ 为公平骰子点数，$\mathcal G$ 只报告奇偶，取
$U=D$、$V=\1_{\{D\geq4\}}$ 以及 $p=q=2$。在奇数块与偶数块上分别有
\[
 \E[UV\mid\mathcal G]=\frac53,\ \frac{10}{3},\qquad
 \E[U^2\mid\mathcal G]=\frac{35}{3},\ \frac{56}{3},\qquad
 \E[V^2\mid\mathcal G]=\frac13,\ \frac23.
\]
所以条件 Cauchy--Schwarz 在两块上分别化为
\[
 \frac53\leq\frac{\sqrt{35}}3,\qquad
 \frac{10}{3}\leq\frac{\sqrt{112}}3;
\]
平方以后就是 $25\leq35$ 与 $100\leq112$。这正是普通 Cauchy--Schwarz 对每个条件分割块的
应用，而不是把两个块的平均混在一起。
\end{example}

端点版本若 $U\in L^1$、$V\in L^\infty$，则
\[
 |\E[UV\mid\mathcal G]|
 \leq\|V\|_\infty\E[|U|\mid\mathcal G].
\]
这是由 $|UV|\leq\|V\|_\infty|U|$ 和条件期望单调性得到的；估计只使用常数界
$\|V\|_\infty$。

\begin{theorem}[$L^2$ 投影性质]\label{theorem:5-2-3-2}
设 $X$ 为实值且 $X\in L^2$。对每个实值 $Y\in L^2(\mathcal G)$，
\begin{equation}\label{eq:5-2-orthogonality}
 \E\bigl[(X-\E[X\mid\mathcal G])Y\bigr]=0.
\end{equation}
因此 $P_{\mathcal G}X=\E[X\mid\mathcal G]$ 是 $L^2(\mathcal G)$ 中唯一最小化
$\E|X-Y|^2$ 的变量。
\end{theorem}

\begin{proof}
记 $P=P_{\mathcal G}X$。由 \eqref{eq:5-2-l2-contraction}，$P\in L^2(\mathcal G)$，故
$X-P\in L^2$。先设 $Y$ 有界。取出已知因子并使用塔式性质，得到
\[
 \begin{aligned}
 \E[(X-P)Y]
 &=\E\bigl[\E[(X-P)Y\mid\mathcal G]\bigr]\\
 &=\E\bigl[Y\E[X-P\mid\mathcal G]\bigr]=0.
 \end{aligned}
\]
对一般 $Y\in L^2(\mathcal G)$，令 $Y_m=(-m)\vee(Y\wedge m)$。\Holder{} 不等式给
\[
 \bigl|\E[(X-P)(Y-Y_m)]\bigr|
 \leq\|X-P\|_2\|Y-Y_m\|_2\longrightarrow0,
\]
所以 \eqref{eq:5-2-orthogonality} 仍成立。

任取 $Y\in L^2(\mathcal G)$，把
$X-Y=(X-P)+(P-Y)$ 展开平方；交叉项由
\eqref{eq:5-2-orthogonality}（取 $P-Y$）为零。因此
\begin{equation}\label{eq:5-2-pythagorean-prediction}
 \|X-Y\|_2^2=\|X-P\|_2^2+\|P-Y\|_2^2.
\end{equation}
右侧第二项非负，且等于零当且仅当 $Y=P$ 于 $L^2$，这同时给出最优性和唯一性。
\end{proof}

\begin{corollary}[总方差公式]\label{corollary:5-2-3-3}
若实值 $X\in L^2$，定义
\[
 \Var(X\mid\mathcal G)
 =\E\bigl[(X-\E[X\mid\mathcal G])^2\mid\mathcal G\bigr].
\]
则
\begin{equation}\label{eq:5-2-total-variance}
 \Var(X)=\E\Var(X\mid\mathcal G)
 +\Var\bigl(\E[X\mid\mathcal G]\bigr).
\end{equation}
\end{corollary}

\begin{proof}
令 $P=\E[X\mid\mathcal G]$。分解
\[
 X-\E X=(X-P)+(P-\E X)
\]
的两项属于 $L^2$，而第二项 $\mathcal G$-可测。由
\eqref{eq:5-2-orthogonality}，它们的内积为零。平方后取期望即得
\[
 \Var(X)=\E(X-P)^2+\E(P-\E X)^2.
\]
塔式性质把第一项写成
$\E\E[(X-P)^2\mid\mathcal G]=\E\Var(X\mid\mathcal G)$，第二项就是
$\Var(P)$。
\end{proof}

\begin{example}[总方差的骰子回算]\label{ex:5-2-3-2}
仍令 $X=D^2$，$\mathcal G$ 只报告奇偶。直接计算
\begin{align*}
 \E X&=\frac{1^2+2^2+3^2+4^2+5^2+6^2}{6}=\frac{91}{6},\\
 \E X^2&=\frac{1^4+2^4+3^4+4^4+5^4+6^4}{6}
 =\frac{2275}{6},\\
 \Var(X)&=\E X^2-(\E X)^2
 =\frac{2275}{6}-\left(\frac{91}{6}\right)^2
 =\frac{5369}{36}.
\end{align*}
奇数块和偶数块内的方差分别为
\[
 \frac13\left[\left(1-\frac{35}{3}\right)^2
 +\left(9-\frac{35}{3}\right)^2
 +\left(25-\frac{35}{3}\right)^2\right]=\frac{896}{9},
\]
\[
 \frac13\left[\left(4-\frac{56}{3}\right)^2
 +\left(16-\frac{56}{3}\right)^2
 +\left(36-\frac{56}{3}\right)^2\right]=\frac{1568}{9}.
\]
两块等概率，故
$\E\Var(X\mid\mathcal G)=1232/9$。条件均值取 $35/3,56/3$，各以概率 $1/2$ 出现，
其均值为 $91/6$，所以
\[
 \Var(\E[X\mid\mathcal G])=\frac12\left(-\frac72\right)^2
 +\frac12\left(\frac72\right)^2=\frac{49}{4}.
\]
最后
\[
 \frac{1232}{9}+\frac{49}{4}=\frac{5369}{36},
\]
与直接方差完全一致。
\end{example}

\begin{example}[联合高斯分布的条件均值]\label{ex:5-2-3-1}
本章称实随机向量 $(X,Y)$ \emph{联合高斯}，若每个实线性组合 $sX+tY$ 或具有一维
高斯分布，或退化为常数。设
\[
 \mu_X=\E X,\quad \mu_Y=\E Y,\quad
 \sigma_X^2=\Var(X),\quad\sigma_Y^2=\Var(Y)>0,
 \quad c=\Cov(X,Y).
\]
先考虑非退化密度模型。令 $\Delta=\sigma_X^2\sigma_Y^2-c^2>0$，并假设联合密度为
\[
 \frac1{2\pi\sqrt\Delta}
 \exp\!\left[-\frac{
 \sigma_Y^2(x-\mu_X)^2-2c(x-\mu_X)(y-\mu_Y)
 +\sigma_X^2(y-\mu_Y)^2}{2\Delta}\right].
\]
关于 $x$ 完成平方：分子中的二次式等于
\[
 \sigma_Y^2\left(x-\mu_X-\frac{c}{\sigma_Y^2}(y-\mu_Y)\right)^2
 +\frac{\Delta}{\sigma_Y^2}(y-\mu_Y)^2.
\]
记
\[
 m(y)=\mu_X+\frac{c}{\sigma_Y^2}(y-\mu_Y)
\]
。把完全平方式代回联合密度，再对 $x$ 作高斯积分，得到
\begin{align*}
 p_Y(y)
 &=\frac{e^{-(y-\mu_Y)^2/(2\sigma_Y^2)}}{2\pi\sqrt\Delta}
   \int_\R
   e^{-\sigma_Y^2(x-m(y))^2/(2\Delta)}\dd x\\
 &=\frac1{\sqrt{2\pi\sigma_Y^2}}
   e^{-(y-\mu_Y)^2/(2\sigma_Y^2)}.
\end{align*}
该密度对每个 $y$ 都严格为正，因而密度比给出条件密度
\[
 p_{X\mid Y=y}(x)
 =\sqrt{\frac{\sigma_Y^2}{2\pi\Delta}}
  \exp\!\left[-\frac{\sigma_Y^2(x-m(y))^2}{2\Delta}\right].
\]
它是均值 $m(y)$、方差 $\Delta/\sigma_Y^2$ 的高斯密度。确实，令
$u=x-m(y)$，则
\begin{align*}
 \int_\R x p_{X\mid Y=y}(x)\dd x
 &=m(y)\int_\R p_{X\mid Y=y}(x)\dd x\\
 &\quad+
 \sqrt{\frac{\sigma_Y^2}{2\pi\Delta}}
 \int_\R u e^{-\sigma_Y^2u^2/(2\Delta)}\dd u
 =m(y),
\end{align*}
因为最后一项是可积奇函数的积分。于是
\begin{equation}\label{eq:5-2-gaussian-regression}
 \E[X\mid Y]
 =\mu_X+\frac{\Cov(X,Y)}{\Var(Y)}(Y-\mu_Y).
\end{equation}
若 $\Delta=0$，则
\[
 \Var\!\left(X-\mu_X-\frac{c}{\sigma_Y^2}(Y-\mu_Y)\right)
 =\sigma_X^2-\frac{c^2}{\sigma_Y^2}=0,
\]
所以括号内变量 a.s. 为零，同一公式仍成立。若 $\Var(Y)=0$，则 $Y$ a.s. 为常数，条件均值为
$\E X$，而 \eqref{eq:5-2-gaussian-regression} 中的逆方差写法不适用。

第~\ref{subsec:5-4-3}~节将证明：协方差正定时，上述密度描述与“所有线性组合均为高斯分布”
的定义一致。因此该计算也给出一般非退化联合高斯变量对的条件均值。一般 $L^2$ 投影只说明
条件均值是所有 $Y$ 的可测函数中的最佳预测，并不能推出它必为仿射函数；仿射性来自高斯
结构。
\end{example}

现在回到换测度。分母可能为零，故贝叶斯公式必须在积分定义中验证，而不能把随机变量形式地约去。

\begin{theorem}[贝叶斯--Radon--Nikodym 公式]\label{theorem:5-2-3-4}
设 $\mathbb Q\ll\Prob$，$Z=\dd\mathbb Q/\dd\Prob$，且
$X\in L^1(\mathbb Q)$。则在
$Z_{\mathcal G}=\E_\Prob[Z\mid\mathcal G]>0$ 处，
\begin{equation}\label{eq:5-2-bayes-rn}
 \E_\mathbb Q[X\mid\mathcal G]
 =\frac{\E_\Prob[ZX\mid\mathcal G]}
 {\E_\Prob[Z\mid\mathcal G]}.
\end{equation}
在 $\{Z_{\mathcal G}=0\}$ 上可任取有限的 $\mathcal G$-可测值；这个集合是
$\mathbb Q$-零集。
\end{theorem}

\begin{proof}
记 $A=Z_{\mathcal G}$、$N=\E_\Prob[ZX\mid\mathcal G]$，并定义
\[
 W=\begin{cases}N/A,&A>0,\\0,&A=0.\end{cases}
\]
因为 $Z\geq0$ 且 $\E_\Prob Z=1$，$A$ 非负、可积并有限 a.s.。在
$E=\{A=0\}\in\mathcal G$ 上，
\[
 \int_EZ\,\dd\Prob=\int_EA\,\dd\Prob=0,
\]
故 $Z=0$ a.s. 于 $E$，并且 $\mathbb Q(E)=0$。基本模长估计给
\[
 |N|\leq\E_\Prob[Z|X|\mid\mathcal G].
\]
右端在 $E$ 上积分为 $\int_EZ|X|\,\dd\Prob=0$，所以 $N=0$ a.s. 于 $E$。

对任意非负 $\mathcal G$-可测函数 $H$，先对简单 $H$ 使用条件期望定义，再由单调收敛推广，得到
\begin{equation}\label{eq:5-2-density-on-information}
 \int HZ\,\dd\Prob=\int HA\,\dd\Prob,
\end{equation}
两边允许为 $+\infty$。因此
\[
 \begin{aligned}
 \int|W|\,\dd\mathbb Q
 &=\int|W|Z\,\dd\Prob
 =\int|W|A\,\dd\Prob\\
 &=\int_{\{A>0\}}|N|\,\dd\Prob
 \leq\int\E_\Prob[Z|X|\mid\mathcal G]\,\dd\Prob
 =\E_\Prob[Z|X|]<\infty.
 \end{aligned}
\]
所以 $W\in L^1(\mathbb Q)$。对 $G\in\mathcal G$，把
\eqref{eq:5-2-density-on-information} 分别用于 $W$ 的正负部或实虚部，得到
\[
 \begin{aligned}
 \int_GW\,\dd\mathbb Q
 &=\int_GWZ\,\dd\Prob
 =\int_GWA\,\dd\Prob\\
 &=\int_GN\,\dd\Prob
 =\int_GZX\,\dd\Prob
 =\int_GX\,\dd\mathbb Q.
 \end{aligned}
\]
这里 $WA=N$ a.s.，包括 $A=0$ 的集合。故 $W$ 满足关于 $\mathbb Q$ 的条件期望定义，
即为 \eqref{eq:5-2-bayes-rn} 的版本。
\end{proof}

\begin{example}[二项--Beta 共轭更新]\label{ex:5-2-3-3}
设 $\alpha,\beta>0$，参数 $\theta\in(0,1)$ 的先验分布为
$\mathrm{Beta}(\alpha,\beta)$，其密度为
\[
 \pi(\theta)=\frac{\theta^{\alpha-1}(1-\theta)^{\beta-1}}
 {B(\alpha,\beta)},\qquad0<\theta<1,
\]
其中
$B(a,b)=\int_0^1t^{a-1}(1-t)^{b-1}\,\dd t$。给定 $\theta$ 后作 $n$ 次独立 Bernoulli
试验，记成功次数为 $K$。似然为
\[
 \Prob(K=k\mid\theta)=\binom nk\theta^k(1-\theta)^{n-k}.
\]
对 $k\in\{0,\ldots,n\}$，参数与观测值的联合密度（对 $\theta$ 使用 Lebesgue 测度，
对 $k$ 使用计数测度）为
\[
 \binom nk\frac{\theta^{\alpha+k-1}(1-\theta)^{\beta+n-k-1}}
 {B(\alpha,\beta)}\1_{(0,1)}(\theta).
\]
因此 $K=k$ 的边缘概率是
\begin{align*}
 \Prob(K=k)
 &=\frac{\binom nk}{B(\alpha,\beta)}
   \int_0^1\theta^{\alpha+k-1}(1-\theta)^{\beta+n-k-1}\dd\theta\\
 &=\binom nk\frac{B(\alpha+k,\beta+n-k)}{B(\alpha,\beta)}.
\end{align*}
由 $\alpha,\beta>0$ 可知这个数严格为正。联合密度除以该边缘概率后，组合数和原归一化常数相消，得到后验密度
\[
 \begin{aligned}
 \pi(\theta\mid K=k)
 &=\frac{\theta^{\alpha+k-1}(1-\theta)^{\beta+n-k-1}}
 {B(\alpha+k,\beta+n-k)},\qquad0<\theta<1.
 \end{aligned}
\]
所以后验为 $\mathrm{Beta}(\alpha+k,\beta+n-k)$。“共轭”在这里表示先验密度乘以似然后，经一次
明确归一化仍落在同一函数族，而不是一条需要背诵的参数口诀。
\end{example}

\begin{figure}[htbp]
\centering
\resizebox{0.98\textwidth}{!}{%
\begin{tikzpicture}[node distance=11mm and 14mm]
  \node[book strong node,minimum width=24mm] (x) {$X\in L^2$};
  \node[book strong node,right=20mm of x,minimum width=31mm] (p) {$P_{\mathcal G}X$\\最佳可观察预测};
  \node[book node,below=of p,minimum width=31mm] (sub) {$L^2(\mathcal G)$};
  \draw[book accent arrow] (x)--node[above,book label,align=center,font=\scriptsize]{条件\\期望}(p);
  \draw[book arrow] (p)--(sub);
  \draw[book dashed arrow] (x.south)--node[below left,book label]{$X-P_{\mathcal G}X\perp L^2(\mathcal G)$}(sub.west);

  \node[book warm node,right=27mm of p,minimum width=25mm] (prior) {先验\\$\pi(\dd\theta)$};
  \node[book node,below=of prior,minimum width=28mm] (like) {似然\\$L(\theta,\dd y)$};
  \node[book strong node,right=24mm of prior,minimum width=29mm] (post) {后验\\$\pi(\dd\theta\mid y)$};
  \draw[book arrow] (prior)--(like);
  \draw[book accent arrow] (prior)--node[above,book label,align=center,font=\scriptsize]{乘似然\\并归一化}(post);
  \draw[book arrow] (like.east)--(post.south);
\end{tikzpicture}
}
\caption{左：条件期望的平方误差几何。右：贝叶斯更新是联合分布相对于观测边缘的测度分解，也可由 Radon--Nikodym 密度比计算。}
\label{fig:5-2-3-projection-bayes}
\end{figure}

条件 Jensen 需要凸性和相应可积性；若 $\varphi(X)$ 不可积，有限值条件期望公式不能直接套用。
$L^2$ 投影解释依赖平方可积，不能无条件移到一般 $L^1$。贝叶斯公式的分母为零时，联合分布不
决定该状态上的点值；定理只在 $\mathbb Q$-a.s. 意义下给出唯一版本。

平方误差分解还可用来比较不同信息层上的预测误差；这类比较所需的结论已由
投影恒等式~\eqref{eq:5-2-pythagorean-prediction} 给出。

\begin{exercise}[条件 Jensen 的等号]
设 $\varphi$ 严格凸，$X,\varphi(X)\in L^1$。证明若
$X=\E[X\mid\mathcal G]$ a.s.，则条件 Jensen 取等号；并在有限分割情形证明反向：若每块有正
概率且取等号，则 $X$ 在每块上 a.s. 为常数。
\end{exercise}

\begin{exercise}[最小均方预测]
令 $X\in L^2$，$Z$ 为标准 Borel 值随机元。证明在所有形如 $g(Z)\in L^2$ 的预测中，
$\E[X\mid Z]$ 唯一最小化均方误差。
\end{exercise}

\begin{exercise}[换测度的塔式检查]
在\cref{theorem:5-2-3-4}的假设下，证明
$\E_\Prob Z_{\mathcal G}=1$，并对 $X=1$ 直接验证贝叶斯--Radon--Nikodym 公式。再说明为什么
$Z_{\mathcal G}=0$ 的集合必为 $\mathbb Q$-零集，却未必为 $\Prob$-零集。
\end{exercise}

条件信息固定以后，随机变量序列仍可用 a.s.、依概率、$L^p$ 与依分布等不同方式趋近。下一节将
画出这些收敛概念之间的蕴含图，并用明确反例标出每条不能反向的箭头。
